% P06.2 scalar regularization two-gate reduction proof, live integration.

\begin{proof}[\Cref{prop:p062-scalar-regularization-two-gate}]
The first condition gives
\[
 \dprimeddprime[h_L]=0,
 \qquad
 \dprimeddprime[-h_L]=-\dprimeddprime[h_L]=0.
\]
Thus both Hessian currents are nonnegative.  The second condition makes
$h_L$ and $-h_L$ CLD-psh, and
\Cref{thm:cld-to-classical-psh} makes both functions classically
plurisubharmonic.  Therefore $h_L$ is classically pluriharmonic.  Scalar
extension invariance for classical pluriharmonicity
\cite[Proposition~7.2(6)]{GR25Classical} descends this conclusion from
$X_L$ to $X$.  The local $\mathbf R$-PL conclusion now follows from
\Cref{thm:supplied-p061-singular-pluriharmonic}.  If $X$ is
pure-dimensional, the original equation $\dprimeddprime[h]=0$ makes the
Hessian currents of both $h$ and $-h$ positive.  The real-PL
positive-current comparison
\Cref{thm:p01-real-pl-positive-current-candidate} makes both signs
semipositive, which is GJR harmonicity.
\end{proof}

\begin{proof}[\Cref{cor:p062-counterexample-two-gate-obstruction}]
Fix an analytic extension $L/K$ and suppose that the first alternative
fails.  Then
\[
 \dprimeddprime[h_L]=\dprimeddprime[-h_L]=0,
\]
so both signs have nonnegative Hessian current.  If both belonged to the
CLD-psh hierarchy, the two conditions of
\Cref{prop:p062-scalar-regularization-two-gate} would hold and $h$ would
be classically pluriharmonic, contrary to the hypothesis.  At least one
sign therefore gives the failure in the second alternative.  Taking
$L=K$ and using the assumed equation $\dprimeddprime[h]=0$ proves the
last assertion.
\end{proof}
